\section{Conclusion}
\label{sec:conclu}
Developing efficient and scalable scene reconstruction methodologies is a fundamental requirement for building closed-loop evaluation and simulation platforms within the autonomous driving industry. In this paper, we introduced ReconDrive, a feed-forward framework for efficient 4D Gaussian Splatting generation specifically tailored for large-scale driving scenes. Unlike existing per-scene optimization methods that require extensive computational resources, our model generates a complete scene representation in a single forward pass. This efficiency is achieved through a static-dynamic composition strategy and segment-wise Gaussian aggregation, which effectively adapt a 3D foundation model via our specialized Gaussian prediction architecture.

Experimental results on our new nuScenes-based benchmark demonstrate that ReconDrive outperforms state-of-the-art per-scene optimization baselines across nearly all evaluation metrics. To the best of our knowledge, this work represents the first instance where a feed-forward approach surpasses optimization-based methods in reconstruction and novel-view synthesis performance through comprehensive experimentation. These findings provide a clear direction and renewed confidence for the development of scalable, generative simulation environments for large-scale autonomous driving.

\paragraph{Limitations and Future Work.} This study represents an initial effort toward achieving large-scale autonomous driving scene reconstruction and novel-view synthesis via feed-forward 3D foundation models. While our results demonstrate significant potential for this paradigm, several challenges remain to be addressed in future research.
\begin{itemize}
    \item Temporal Representation of Non-Rigid Motion: ReconDrive currently utilizes an explicit 4D Gaussian Splatting representation based on linear motion estimation. While efficient, this formulation may struggle to accurately represent complex non-rigid deformations or highly non-linear object trajectories. Future work could involve integrating more expressive temporal kernels or learnable deformation fields to better capture these dynamics.
    \item Temporal Consistency and Redundancy in Aggregation: Our current post-processing strategy for multi-frame aggregation relies on pixel-wise outputs that may lead to Gaussian redundancy and suboptimal handling of occluded regions. Transitioning toward a more integrated temporal fusion mechanism within the latent space could reduce redundancy and improve the structural integrity of occluded surfaces.
    \item Precision in Dynamic Object Extraction: The current pipeline relies on SAM2 for object segmentation, which can occasionally lead to inaccurate boundaries or missed detections. Furthermore, displacing dynamic objects directly can leave disocclusion artifacts in the background. Investigating joint inpainting techniques or end-to-end motion modeling could help mitigate these "background holes" and improve boundary fidelity.
    \item Computational Efficiency and Throughput: Although our feed-forward approach is orders of magnitude faster than optimization-based methods, there remains substantial room for architectural optimization. Researching lightweight foundation backbones and more efficient Gaussian sampling strategies will be critical for achieving real-time performance in edge-computing environments.
    \item Scalability and Generalization: Expanding the training scope to encompass larger, more diverse datasets is essential. Evaluating the model across a broader range of geographic locations and extreme weather conditions will be a primary focus to ensure the robustness and generalizability of the feed-forward paradigm. Broadening data diversity to cover different driving behaviors may further enhance robustness across varied real-world scenarios~\cite{hao2025styledrive}.
\end{itemize}